% 文件名：简明LaTeX教程.tex
% 编译命令：xelatex 简明LaTeX教程.tex
% 再编译一次：xelatex 简明LaTeX教程.tex
% 目的：第一次生成目录信息，第二次目录页码和条目才正确

\documentclass[12pt,a4paper]{ctexart}

% ===== 颜色与链接 =====
\usepackage{xcolor}
\usepackage[colorlinks=true, linkcolor=blue, urlcolor=blue, citecolor=blue]{hyperref}
% 或者无边框：
% \usepackage[hidelinks]{hyperref}

% ===== 表格 =====
\usepackage{array}
\usepackage{booktabs}
\usepackage{tabularx}
\usepackage{multirow}

% ===== 数学 =====
\usepackage{amsmath}
\usepackage{amssymb}

% ===== 图片 =====
\usepackage{graphicx}

% ===== 页面 =====
\usepackage{geometry}
\geometry{margin=2.5cm}

% ===== 编号设置 =====
\setcounter{tocdepth}{1}
\setcounter{secnumdepth}{1}
\counterwithin{table}{section}   % 表 1.1
% \counterwithin{figure}{section} % 图 1.1，需要时取消注释

% ===== 文档信息 =====
\title{简明LaTeX教程}
\author{}
\date{}

\begin{document}

% 避免 Overfull \hbox 警告：允许 TeX 在实在断不开的长行上多拉伸一点词间距。
% 只在"无法正常断行"的那几行生效，不会影响其他段落的正常排版。
\setlength{\emergencystretch}{3em}

\maketitle
\clearpage

\tableofcontents
\clearpage

% ==================== 基础部分 ====================

\section{源代码结构}
一个简单的LaTex程序主要包括导言区和正文区，导言区主要设置文档类和导入宏，正文区是主要的写作区域。一个具体的LaTex程序如下：

\texttt{\textbackslash documentclass\{cetxart\}}\texttt{\% 文档类}

\texttt{\textbackslash usepackage\{ctex\}}\texttt{\% 宏包}

\texttt{\textbackslash begin\{document\}}

\texttt{\% 正文区}

Hello World!

\texttt{\textbackslash end\{document\}}

将它们编译运行后的结果是：

Hello World!

下面将讲解这段程序。

LaTeX 源文件的扩展名是 \texttt{.tex}。最简单的文档类声明格式如下：

\begin{center}
  \texttt{\textbackslash documentclass[<options>]\{<classname>\}}
\end{center}

其中 \texttt{<classname>} 为文档类名称，常用文档类有
\texttt{article}、\texttt{report}、\texttt{book}、\texttt{proc}、
\texttt{slides}、\texttt{minimal}，以及支持中文排版的
\texttt{ctexart}、\texttt{ctexrep}、\texttt{ctexbook} 等，如表~\ref{tab:classname}

\begin{table}[htbp]
  \centering
  \caption{LaTeX 常用文档类一览}
  \label{tab:classname}
  \begin{tabularx}{\textwidth}{@{}l l X@{}}
    \toprule
    文档类 & 中文含义 & 用途与主要特点 \\
    \midrule
    article & 文章 & 用于短文、论文、报告；无 \texttt{\textbackslash chapter}，章节命令为 \texttt{\textbackslash section}、\texttt{\textbackslash subsection} 等。 \\
    report & 报告 & 用于较长报告、学位论文；有 \texttt{\textbackslash chapter}，可含摘要、附录等。 \\
    book & 书 & 用于书籍排版；有 \texttt{\textbackslash part}、\texttt{\textbackslash chapter}；通常双面排版。 \\
    proc & 会议论文集 & 用于会议论文；基于 article，提供会议论文常用格式。 \\
    slides & 幻灯片 & 用于制作投影幻灯片；较老旧，现代推荐用 beamer。 \\
    minimal & 最小文档 & 仅用于最小测试；几乎不提供标题、章节等排版功能。 \\
    ctexart & 中文文章 & 支持中文的 article；自动处理中文字体、标点、断行等。 \\
    ctexrep & 中文报告 & 支持中文的 report。 \\
    ctexbook & 中文书籍 & 支持中文的 book。 \\
    \bottomrule
  \end{tabularx}
\end{table}

文档类下边需要有：

 \texttt{\textbackslash begin\{document\}}

 \texttt{\% 正文区}

  Hello World!

 \texttt{\textbackslash end\{document\}}

  在\texttt{\textbackslash begin\{document\}}和 \texttt{\textbackslash end\{document\}}之间的区域就是正文区，我们写的文字所有内容都需要放在正文区内。
  \texttt{\% 文档类}{\% 正文区}是注释，在LaTex的源文件中\texttt{\%}后的东西是注释，注释是对文档的解释说明，编译器自动忽略，不会编译到文件当中。

  在LaTex中我们无法直接输入\texttt{\textbackslash 、 \% 、 \{ 、 \} 等内容}，需要使用转义字符来表示，这里将常见的转义字符总结为表~\ref{tab:escape-chars}。

  \begin{table}[htbp]
  \centering
  \caption{LaTeX 常见转义字符}
  \label{tab:escape-chars}
  \begin{tabularx}{\textwidth}{@{}c c X@{}}
    \toprule
    字符 & LaTeX 输入 & 说明 \\
    \midrule
     \texttt{\textbackslash} & \texttt{\textbackslash textbackslash} & 命令开始符，文本中需用 \texttt{\textbackslash textbackslash} 输出。 \\
     \texttt{\%} & \texttt{\textbackslash \%} & 注释符，需转义。 \\
     \texttt{\{} & \texttt{\textbackslash \{} & 分组开始，需转义。 \\
     \texttt{\}} & \texttt{\textbackslash \}} & 分组结束，需转义。 \\
     \texttt{\$} & \texttt{\textbackslash \$} & 数学模式开关，需转义。 \\
     \texttt{\&} & \texttt{\textbackslash \&} & 表格列分隔符，需转义。 \\
     \texttt{\#} & \texttt{\textbackslash \#} & 宏参数序号，需转义。 \\
     \texttt{\_} & \texttt{\textbackslash \_} & 下标命令，需转义。 \\
     \texttt{\textasciicircum} & \texttt{\textbackslash textasciicircum} & 上标命令，文本中需用 \texttt{\textbackslash textasciicircum}。 \\
     \texttt{\textasciitilde} & \texttt{\textbackslash textasciitilde} & 非断行空格，文本中需用 \texttt{\textbackslash textasciitilde}。 \\
    \bottomrule
  \end{tabularx}
\end{table}

以\texttt{\textbackslash begin}为开头，以 \texttt{\textbackslash end} 为结尾的区域叫作环境，在这个区域之前的区域为导言区。

  所有 \texttt{\textbackslash} 后带单词的部分称为命令，常见的命令有
  \texttt{\textbackslash begin}、\texttt{\textbackslash end}、
  \texttt{\textbackslash document}、\texttt{\textbackslash usepackage}、
  \texttt{\textbackslash clearpage} 等。
  其中\texttt{\textbackslash usepackage}的作用是引用宏，当文章需要使用中文，但不能使用支持中文的文档类时，可以使用
  \texttt{\textbackslash usepackage(ctex)}来导入中文宏，这样就可以使用中文写作了。

\section{标题、作者、日期}
这一节介绍如何给文章添加标题、作者、日期。

通常，可以在导言区完成文章标题、作者、日期的设置，具体方式为：

使用\texttt{\textbackslash title \{标题\}}命令为文章添加标题。大括号中写入文章的标题。

\texttt{\textbackslash author \{作者\}}命令同理，用于为文章添加作者。作者可以通过\texttt{\textbackslash and}命令添加多个作者，在作者后添加\texttt{\textbackslash thanks}命令会为作者添加附加信息。

有时文章需要添加写作的日期，使用\texttt{\textbackslash date\{日期\}}命令可以完成这个操作。如果日期填入\texttt{\textbackslash today}这个命令，程序会显示编译文档的日期。

这些设置是在导言区完成的，直接编译不会在pdf中显示出来，需要在正文区中写入\texttt{\textbackslash maketitle}命令，这样程序就可以显示这些设置了，\texttt{\textbackslash maketitle}的位置就是这些设置显示的位置。

标题、作者、日期这些基础样式在文档类中已经设置好了，不同文档类的排版是不一样的。


\section{换行、分段、空格}
在LaTex中空格不等于换行。比如：

include<stdio.h>

int main(void)

\{

\texttt{  printf("Hello World!");}

\}

实际编译后为：

include<stdio.h>
int main(void)
{
printf("Hello World!");
}

在LaTex中我们可以使用\texttt{\textbackslash \textbackslash}、\texttt{\textbackslash newline}来表示换行，换行后首行不会缩进。

直接空一行或多行是用来表示另起一段的常用方式，也可以使用\texttt{\textbackslash par}来表示另起一段，分段首行会进行缩进。

\texttt{\textbackslash hspace\{n em(pt/cm/in/ex)\}}\allowbreak 命令用于打空格，大括号内表示打出n个字符(磅/厘米/英寸/字符x) 大小的空格。

\texttt{\textbackslash baselineskip} 表示当前字号下的一行高度，所以 \texttt{\textbackslash vspace*\{baselineskip\}} 就是空出一行。
\texttt{\textbackslash vspace*} 带星号，表示即使处在段落开头也强制插入空白。


\section{居中、取消缩进、行距}
使用

\texttt{\textbackslash begin\{center\}}

\texttt{\textbackslash end\{center\}}

设置居中环境，这个环境中的所有行都是居中显示的。

同理，也可以设置flushright、flushleft居右和具左环境。

\vspace*{\baselineskip}

\texttt{\textbackslash noindent} 是取消缩进命令，哪个位置要取消缩进就写在哪个位置；
在导言区写入\texttt{\textbackslash parindent = 0pt}通过设置缩进为0，同样能够做到取消缩进；
\texttt{\textbackslash parindent = 0pt}还可以写成\texttt{\textbackslash setlength\{parindent\}\{0pt\}}，它们具有同样效果。

\vspace*{\baselineskip}

在导言区通过\texttt{\textbackslash linespread\{行距倍数\}}命令能够设置全篇的行距倍数；在正文区使用\texttt{\textbackslash linespread\{行距倍数\}}\texttt{\textbackslash selestfont}会修改以这个命令为起始，它之后所有正文的行距；
如果用大括号将\texttt{\textbackslash linespread\{行距倍数\}}\texttt{\textbackslash selestfont}和正文一起括起来并分段，即\{\texttt{\textbackslash linespread\{行距倍数\}}\texttt{\textbackslash selestfont}正文\textbackslash par\}，就可以控制大括号内的正文部分行距。

在导言区使用\texttt{\textbackslash setstretch\{行距倍数\}}也同样可以设置全局行距。

还可以在导言区使用命令\texttt{\textbackslash udepackage\{setspace\}}导入宏包，这个宏包里定义了\texttt{\textbackslash singlespacing}、\texttt{\textbackslash onehalfspacing}和\texttt{\textbackslash doublespacing}单倍行距、1.5倍行距和双倍行距命令，这个行距与其他行距不同，指的是字号而不是基础行距。

通过设置\texttt{\textbackslash begin\{spacing\}\{行距倍数\}}和\texttt{\textbackslash end\{spacing\}}，即spacing环境，也能够修改环境内的行距。

不同的软件中，行距的定义可能不同。在world和LaTex两个排版工具中，行距倍数指的是基础行距的倍数。
在LaTex中，基础行距的定义是字号的1.2倍，在word中是字号的1.3倍需要注意转化。


\section{字号与字体}
当没有任何排版设置时，中文排版的字号大小是11磅，这个字号被称为正文字号，就是文章主体字号的大小。

在导入文档类添加\texttt{\textbackslash documentclass[10/11/12pt]\{文档类\}}，可以设置整个文章主题的字号的大小。

通过字号命令可以修改自字号以下环境/花括号括中的正文字号；字号命令是相对命令，实际大小取决于文档类的正文字号。
表~\ref{tab:font-size-pt} 给出常用字号命令在 10pt、11pt、12pt 文档中的实际磅值。

\begin{table}[htbp]
  \centering
  \caption{不同正文字号下字号命令的实际磅值}
  \label{tab:font-size-pt}
  \begin{tabular*}{\textwidth}{@{\extracolsep{\fill}}l l c c c@{}}
    \toprule
    \multirow{2}{*}{样本} & \multirow{2}{*}{命令}
      & \multicolumn{3}{c}{正文字号} \\
    \cmidrule(lr){3-5}
      & & 10pt & 11pt & 12pt \\
    \midrule
    {\tiny sample}         & \texttt{\textbackslash tiny}         & 5pt  & 6pt  & 6pt  \\
    {\scriptsize sample}   & \texttt{\textbackslash scriptsize}   & 7pt  & 8pt  & 8pt  \\
    {\footnotesize sample} & \texttt{\textbackslash footnotesize} & 8pt  & 9pt  & 10pt \\
    {\small sample}        & \texttt{\textbackslash small}        & 9pt  & 10pt & 11pt \\
    {\normalsize sample}   & \texttt{\textbackslash normalsize}   & 10pt & 11pt & 12pt \\
    {\large sample}        & \texttt{\textbackslash large}        & 12pt & 12pt & 14pt \\
    {\Large sample}        & \texttt{\textbackslash Large}        & 14pt & 14pt & 17pt \\
    {\LARGE sample}        & \texttt{\textbackslash LARGE}        & 17pt & 17pt & 20pt \\
    {\huge sample}         & \texttt{\textbackslash huge}         & 20pt & 20pt & 25pt \\
    {\Huge sample}         & \texttt{\textbackslash Huge}         & 25pt & 25pt & 25pt \\
    \bottomrule
  \end{tabular*}
\end{table}

使用\texttt{\textbackslash fontsize\{字体大小(pt)\}\{基础行距(pt)\}}，同样能够修改环境或花括号中以它为起点，往后的正文字号。

导入中文宏包（如 \texttt{ctex}）后，可用 \texttt{\textbackslash zihao\{\}} 命令按中文号数设置字号。
小号字在号数前加负号，例如 \texttt{\textbackslash zihao\{-4\}} 是小四。
各号数与 bp、pt 的对应关系见表~\ref{tab:zihao}。

\begin{table}[htbp]
  \centering
  \caption{\texttt{\textbackslash zihao} 字号对照表}
  \label{tab:zihao}
  \begin{tabularx}{\textwidth}{@{}l l c c X@{}}
    \toprule
    字号 & \texttt{\textbackslash zihao} 参数 & 大小（bp） & 大小（pt） & 示例 \\
    \midrule
    初号   & \texttt{\{0\}}   & 42 bp   & 42.1575 pt & {\zihao{0} 字号} \\
    小初   & \texttt{\{-0\}}  & 36 bp   & 36.1350 pt & {\zihao{-0} 字号} \\
    一号   & \texttt{\{1\}}   & 26 bp   & 26.0975 pt & {\zihao{1} 字号} \\
    小一   & \texttt{\{-1\}}  & 24 bp   & 24.0900 pt & {\zihao{-1} 字号} \\
    二号   & \texttt{\{2\}}   & 22 bp   & 22.0825 pt & {\zihao{2} 字号} \\
    小二   & \texttt{\{-2\}}  & 18 bp   & 18.0675 pt & {\zihao{-2} 字号} \\
    三号   & \texttt{\{3\}}   & 16 bp   & 16.0600 pt & {\zihao{3} 字号} \\
    小三   & \texttt{\{-3\}}  & 15 bp   & 15.0563 pt & {\zihao{-3} 字号} \\
    四号   & \texttt{\{4\}}   & 14 bp   & 14.0525 pt & {\zihao{4} 字号} \\
    小四   & \texttt{\{-4\}}  & 12 bp   & 12.0450 pt & {\zihao{-4} 字号} \\
    五号   & \texttt{\{5\}}   & 10.5 bp & 10.5394 pt & {\zihao{5} 字号} \\
    小五   & \texttt{\{-5\}}  & 9 bp    & 9.0338 pt  & {\zihao{-5} 字号} \\
    六号   & \texttt{\{6\}}   & 7.5 bp  & 7.5281 pt  & {\zihao{6} 字号} \\
    小六   & \texttt{\{-6\}}  & 6.5 bp  & 6.5244 pt  & {\zihao{-6} 字号} \\
    七号   & \texttt{\{7\}}   & 5.5 bp  & 5.5206 pt  & {\zihao{7} 字号} \\
    八号   & \texttt{\{8\}}   & 5 bp    & 5.0188 pt  & {\zihao{8} 字号} \\
    \bottomrule
  \end{tabularx}
\end{table}

\section{下划线与颜色}
下划线用 \verb|\underline|，颜色需要 \verb|xcolor| 宏包，
用 \verb|\textcolor| 设置前景色，\verb|\colorbox| 设置背景色。

\section{纸张与边距}
纸张在 \verb|\documentclass| 里设置，如 \verb|a4paper|；
页边距用 \verb|geometry| 宏包，例如 \verb|\geometry{margin=2.5cm}|。

\section{分页与分栏}
分页用 \verb|\newpage| 或 \verb|\clearpage|；
分栏需要 \verb|multicol| 宏包，用 \verb|multicols| 环境。

\section{摘要摘录脚注边注}
摘要用 \verb|abstract| 环境，引用块用 \verb|quote| 或 \verb|quotation|，
脚注用 \verb|\footnote|，边注用 \verb|\marginpar|。

\section{盒子}
常用盒子命令有 \verb|\mbox|、\verb|\fbox|、\verb|\parbox|，
更复杂排版可用 \verb|minipage| 环境。

\section{自定义命令}
用 \verb|\newcommand| 定义新命令，用 \verb|\renewcommand| 修改已有命令。

% ==================== 图片部分 ====================

\section{插入图片}
需要 \verb|graphicx| 宏包，用 \verb|\includegraphics| 插入图片，
可指定 \verb|width| 或 \verb|height| 控制大小。

\section{浮动体与图片排版}
图片通常放在 \verb|figure| 浮动体里，用 \verb|\caption| 加标题，
\verb|[htbp]| 控制浮动位置。

\section{图片引用}
用 \verb|\label| 给图片加标签，用 \verb|\ref| 引用图片编号。

% ==================== 列表部分 ====================

\section{有序列表无序列表描述环境}
无序列表用 \verb|itemize|，有序列表用 \verb|enumerate|，
描述列表用 \verb|description|。

\section{一些空格与填充}
水平空格用 \verb|\quad|、\verb|\qquad|、\verb|\hspace|，
垂直空格用 \verb|\vspace|，弹性填充用 \verb|\hfill|、\verb|\vfill|。

% ==================== 表格部分 ====================

\section{简单表格-tabular}
用 \verb|tabular| 环境，\verb|&| 分列，\verb|\\| 换行，\verb|\hline| 画横线。

\section{表格浮动环境与表格引用}
表格通常放在 \verb|table| 浮动体里，用 \verb|\caption| 加标题，
用 \verb|\label| 加标签，用 \verb|\ref| 引用。

\section{一些表格样式}
调整列宽用 \verb|p{}|、\verb|m{}|、\verb|b{}|，
调整行高用 \verb|\arraystretch|。

\section{三线表}
用 \verb|booktabs| 宏包，\verb|\toprule|、\verb|\midrule|、\verb|\bottomrule| 绘制三线。

\section{自动变宽表格-tabularx}
用 \verb|tabularx| 宏包，\verb|X| 列会自动分配剩余宽度。

\section{形式表格-tabularray}
用 \verb|tabularray| 宏包的 \verb|tblr| 环境，语法更现代。

\section{表格-换行与Q格式}
\verb|tabularray| 支持在单元格内换行，用 \verb|Q| 格式控制列样式。

\section{表格-变宽}
\verb|tabularray| 用 \verb|X| 列自动变宽。

\section{表格-行列设置分离}
\verb|tabularray| 可以把行设置和列设置分开写，更清晰。

\section{表格-框线样式颜色}
\verb|tabularray| 支持自定义框线宽度、颜色、样式。

\section{表格-按行列绘制边框线}
可以只给某几行或某几列画边框线。

\section{表格-多边框线}
可以画双线、粗线、虚线等多种边框线。

\section{表格-自由绘制边框线}
可以按坐标自由绘制任意位置的边框线。

\section{表格-设置行列样式}
可以为整行或整列统一设置字体、颜色、对齐等样式。

\section{表格-单元格合并于样式}
用 \verb|\multicolumn|、\verb|\multirow| 合并单元格，
\verb|tabularray| 也支持 \verb|cell| 命令合并和设置样式。

\section{tabularray于三线表}
\verb|tabularray| 可以很方便地画三线表。

\section{长表格-longtblr}
用 \verb|longtblr| 环境制作跨页长表格。

\section{横向表格（宽表）}
用 \verb|sidewaystable| 环境或 \verb|pdflscape| 宏包把宽表旋转。

\section{关于表格的一些补充}
补充说明表格的常见问题，如列宽、行高、跨页、对齐等。

% ==================== 公式部分 ====================

\section{数字模式与公式引用}
行内公式用 \verb|$...$|，独立公式用 \verb|\[...\]| 或 \verb|equation| 环境，
用 \verb|\label| 和 \verb|\ref| 引用公式编号。

\section{分式与根式}
分式用 \verb|\frac{a}{b}|，根式用 \verb|\sqrt{x}| 或 \verb|\sqrt[n]{x}|。

\section{公式-上下标}
上标用 \verb|^|，下标用 \verb|_|，同时上下标用 \verb|x_i^2|。

\section{公式-上下线}
上划线用 \verb|\overline|，下划线用 \verb|\underline|，
上下括号用 \verb|\overbrace|、\verb|\underbrace|。

\section{公式-求和与公式样式调整}
求和用 \verb|\sum|，可加 \verb|\limits| 控制上下限位置，
用 \verb|\displaystyle| 调整公式样式。

\section{公式-积分}
积分用 \verb|\int|，多重积分用 \verb|\iint|、\verb|\iiint|。

\section{公式-矩阵}
矩阵用 \verb|matrix|、\verb|pmatrix|、\verb|bmatrix| 等环境。

\section{公式-多行公式}
多行公式用 \verb|align| 或 \verb|gather| 环境，用 \verb|&| 对齐。

\section{公式-分支公式}
分支公式用 \verb|cases| 环境。

\section{公式-定理环境与cleveref引用}
定理环境用 \verb|amsthm| 宏包，引用可用 \verb|cleveref| 宏包自动识别类型。

% ==================== 文档结构部分 ====================

\section{章节标题与目录}
用 \verb|\section|、\verb|\subsection| 等命令生成标题和目录，
用 \verb|\tableofcontents| 生成目录。

\section{文档结构-前言正文附录后记}
用 \verb|\frontmatter|、\verb|\mainmatter|、\verb|\backmatter| 划分文档结构，
用 \verb|\appendix| 开始附录。

\section{文档拆分}
用 \verb|\include| 或 \verb|\input| 把大文档拆成多个文件。

\section{页眉页脚}
用 \verb|fancyhdr| 宏包自定义页眉页脚。

\section{参考文献}
用 \verb|thebibliography| 环境或 \verb|bibtex|、\verb|biblatex| 管理参考文献。

\section{索引}
用 \verb|makeidx| 或 \verb|imakeidx| 宏包生成索引。

% ==================== TikZ 部分 ====================

\section{TikZ-绘图 直线 坐标}
用 \verb|tikz| 宏包，\verb|\draw (0,0) -- (2,1);| 画一条直线。

\section{TikZ-基本图形绘制}
可以画圆、矩形、椭圆、多边形等基本图形。

\section{TikZ-线条样式}
可以设置线条粗细、颜色、虚实、箭头等样式。

\section{TikZ-scope环境、样式复用}
用 \verb|scope| 环境分组，用 \verb|tikzset| 定义可复用样式。

\section{TikZ-形变}
用 \verb|scale|、\verb|rotate|、\verb|xshift|、\verb|yshift| 等做形变。

\section{TikZ-绘制与填充}
用 \verb|\draw| 绘制，用 \verb|\fill| 填充，用 \verb|filldraw| 同时绘制和填充。

\section{TikZ-文字节点}
用 \verb|\node| 添加文字节点，可设置位置、对齐、颜色等。

\section{TikZ-循环绘制与坐标轴}
用 \verb|\foreach| 循环绘制，用 \verb|axis| 环境画坐标轴。

\section{TikZ-绘制函数图像}
用 \verb|pgfplots| 宏包绘制函数图像。

% ==================== 其他部分 ====================

\section{初识beamer幻灯片}
\verb|beamer| 是制作幻灯片的文档类，用 \verb|frame| 环境表示每一页。

\section{CTeX补充内容}
CTeX 提供中文排版支持，包括字体、标点、断行、章节标题中文化等。

\section{一些目录样式}
可以用 \verb|titletoc|、\verb|tocloft| 等宏包自定义目录样式。

\section{为页面设置背景图片}
用 \verb|background| 或 \verb|eso-pic| 宏包给页面加背景图片。

\section{cutwin文字环绕排版}
用 \verb|cutwin| 宏包实现文字环绕图片的排版效果。

\end{document}
